\section{Adam Response, Memory and Curvature}
\label{appendix:adam-memory}

This appendix relates the raw-gradient response to a preconditioned-curvature
proxy. The identities specify the assumptions needed for that interpretation;
they are not an empirical calibration on the neural trajectories.

For a full objective that is twice continuously differentiable along
the update, with $g_t\ne0$ and $\eta>0$, write
\begin{equation}
P_t:=\operatorname{diag}
\left((\sqrt{\widehat v_t}+\epsilon)^{-1}\right),
\qquad d_t=P_t\widehat m_t,
\end{equation}
and define the starting Hessian $H_t=\nabla^2\mathcal L(\theta_{t-1})$
and its average along the update,
\begin{equation}
\overline H_t:=
\int_0^1\nabla^2\mathcal L(\theta_{t-1}-u\eta d_t)\,du.
\end{equation}
The gradient integral gives
$g_{t+1}-g_t=-\eta\overline H_td_t$, hence
\begin{equation}
\frac{z_{t+1}}{\eta}
=
\frac{g_t^\top\overline H_tP_t\widehat m_t}{\|g_t\|^2}.
\label{eq:adam-response-integral}
\end{equation}
Thus the response is a mixed curvature projection involving the current
slope and realized Adam direction. The classical-Hessian identity requires
this path regularity; it is not automatically valid across nonsmooth
activation boundaries. The projected response itself remains computable
where the recorded gradients exist.

With our indexing, the frozen-time directional proxy of
Bai et al.~\cite{bai2026adaptive} is
\begin{equation}
\begin{aligned}
\lambda_{B,t}
&:=c_t\frac{g_t^\top P_tH_tg_t}{\|g_t\|^2},\\
c_t&:=\frac{1-\beta_1}{(1+\beta_1)(1-\beta_1^t)}.
\end{aligned}
\label{eq:adam-bai-reference}
\end{equation}
Let $e_t:=\widehat m_t-c_tg_t$. Hessian symmetry yields the exact
discrepancy decomposition
\begin{equation}
\begin{aligned}
\frac{z_{t+1}}{\eta}-\lambda_{B,t}
={}&\frac{g_t^\top H_tP_te_t}{\|g_t\|^2}\\
&+\frac{g_t^\top(\overline H_t-H_t)d_t}{\|g_t\|^2}.
\end{aligned}
\label{eq:adam-response-discrepancy}
\end{equation}
The first term measures the effect of departure from the proportional
first-moment response $c_tg_t$. The second measures variation of curvature
along the finite update. This separates the difference between two
diagnostics; causal attribution of spike formation requires additional
controlled experiments. A single HVP $h_t=H_tg_t$ computes both
$\lambda_{B,t}=c_t(P_tg_t)^\top h_t/\|g_t\|^2$ and the first term
$(P_te_t)^\top h_t/\|g_t\|^2$. Their difference from $z_{t+1}/\eta$
has not yet been measured on these neural trajectories.

The first-moment departure is directly connected to adjacent gradient
sums. Let $a=(1-\beta_1)/(1+\beta_1)$ and $r_t^{(m)}=m_t-ag_t$. Substitution
into Eq.~\eqref{eq:adam-moments} gives
\begin{equation}
r_t^{(m)}=\beta_1r_{t-1}^{(m)}+\beta_1a(g_t+g_{t-1}),
\qquad t\ge2,
\label{eq:adam-memory-residual}
\end{equation}
with $e_t=r_t^{(m)}/(1-\beta_1^t)$. The adjacent sums used in the GD pairing
analysis thus drive departure from the alternating first-moment response.
Settled constant-amplitude alternation makes $e_t$ small. The approximation
$z_{t+1}\approx\eta\lambda_{B,t}$ additionally requires both weighted
correction terms in Eq.~\eqref{eq:adam-response-discrepancy} to be small;
the preconditioner can magnify a small memory residual.

For exact alternation starting at $m_0=0$, the finite-time expression is
$m_t=a[1-(-\beta_1)^t]g_t$; a later alternating segment additionally
inherits its initial memory. Near-opposition alone therefore does not
establish the proportionality when gradient magnitudes or directions
change rapidly.

If the Hessian is $M$-Lipschitz along the step, the discrepancy obeys
\begin{equation}
\begin{aligned}
|z_{t+1}-\eta\lambda_{B,t}|
\leq{}&
\eta\frac{\|H_tg_t\|\,\|P_te_t\|}{\|g_t\|^2}\\
&+\frac{M\eta^2\|d_t\|^2}{2\|g_t\|}.
\end{aligned}
\label{eq:adam-response-error-bound}
\end{equation}
The bound follows from Cauchy--Schwarz and
$\|\overline H_t-H_t\|\leq M\eta\|d_t\|/2$.
It characterizes diagnostic approximation under the stated regularity
conditions; it does not assert a general theorem for the occurrence
or recovery of Adam spikes.
